% SPDX-FileCopyrightText: 2026 Yiyu Zhou <yiyu@yiyuzhou.io>
%
% SPDX-License-Identifier: GPL-3.0-or-later

\chapter{Sequences}
\label{ch:sequences}

So far, each cue has needed its own GO\@.  A cue can also fire the next cue
by itself, so that a single GO plays a whole sequence.  Three settings on the
\ui{Basics} tab control the timing:

\begin{description}
\item[Pre-wait] Seconds between the cue firing and its action starting.
\item[Post-wait] Seconds between the action starting and the next cue
  firing, for a cue that continues to the next.
\item[Then] What happens to the next cue: \ui{wait for GO}, the default,
  leaves it to the operator; \ui{auto-continue} fires it when the post-wait
  ends; and \ui{auto-follow} fires it when this cue's action ends.
\end{description}

The \texttt{Pre}, \texttt{Post}, and \texttt{Then} columns of the cue list
show the same settings.

\section{A sequence in the storm}

Press \key{Control+o} and open \texttt{storm.cinq.toml} again.

\begin{enumerate}
\item Select cue~1, press \ui{Standby selected}, and press \ui{GO} to restart
  the preshow music.  Cue~3 is now on standby, and the transport summarizes
  what its GO fires: three cues, the second 4.0 seconds and the third 5.5
  seconds after GO, as Figure~\ref{fig:playing} in
  Chapter~\ref{ch:running} shows.
\item Press \ui{GO}\@.  Cue~3 fades the music out over six seconds.  Four
  seconds after GO, the distant thunder starts, followed 1.5 seconds later by
  the rain, as in Figure~\ref{fig:running}.
\end{enumerate}

\begin{figure}[H]
  \centering
  \includegraphics[width=\textwidth]{screenshots/running.png}
  \caption{Five seconds after GO on cue 3, which fades out the preshow
    music, while the thunder plays and the rain is about to start}
  \label{fig:running}
\end{figure}

Cue~3 auto-continues to cue~4 after a post-wait of 4 seconds, and cue~4
continues to cue~5 after 1.5 seconds.  Cue~5 waits for GO, so the sequence
ends there.  While a cue waits to start, the sidebar counts down to it.

\section{Pre-wait and post-wait}

A pre-wait can produce the same timing, but it belongs to the cue that waits
rather than to the cue before it.

\begin{enumerate}
\item Press \ui{Panic}.  Set the \ui{Post-wait} of cue~3 to 0 and the
  \ui{Pre-wait} of cue~4 to 4.
\item Put cue~1 on standby, press \ui{GO} for the music, and press \ui{GO}
  again for cue~3.  The sequence sounds the same: cue~4 now fires at once and
  waits four seconds before playing.
\item Put cue~4 on standby and press \ui{GO}\@.  The thunder now waits four
  seconds; before the change, it would have started immediately.  A pre-wait
  delays its own cue's action however the cue is fired, while a post-wait
  delays only the next cue.
\item Restore the original values: a post-wait of 4 on cue~3 and a pre-wait
  of 0 on cue~4.
\end{enumerate}

\section{Auto-follow}

Auto-follow fires the next cue when the current cue's action ends, such as
when a sound finishes playing, so it needs no post-wait.  An action that
never ends, such as a looping sound or a picture, never triggers an
auto-follow.

\section{Wait and memo cues}

A wait cue lasts for its \ui{Duration} and does nothing else; it paces a
sequence.  A memo cue does nothing at all.  It serves as a note in the list
or as the start of a sequence, like cue~1 of the storm show, whose notes tell
the operator what the music is for.

\section*{Exercise}

You can now make one GO fire a whole sequence.  In your own show, move the
thunder cue from Chapter~\ref{ch:cues} above the rain with
\key{Alt+ArrowUp}, and set its \ui{Then} to \ui{auto-follow}, so that one GO
plays the thunder and then the rain.
